% !TeX root = ../navier-stokes-manuscript.tex
\subsection{The Tower opens toward more spatial regularity}\label{sub:ThePressureCompleteMixedDerivativeTower}

The proposed failure drove the vortex toward smaller scales because spatial
control was supposed to disappear there.  Finite time then drove the same
history upward through its temporal responses.  But the equation that
generates each response contains \(\nu\Delta\) of the response below it.  The
attempt to lose spatial regularity has opened an algebraic direction toward
more of it.  We now have to find out what that direction actually pays for.

The critical height puts the temporal and spatial movements on the same scale.
Under
\eqref{eq:BodyNavierStokesScaling},
\[
  H_\lambda(t)=H(\lambda^2t),
  \qquad
  H_\lambda^{(n)}(t)
  =\lambda^{2n}H^{(n)}(\lambda^2t).
\]
The spatial energy at height \(n+\tfrac12\) carries exactly that weight:
\[
  E_{n+1/2,\lambda}(t)
  =\lambda^{2n}E_{n+1/2}(\lambda^2t).
\]
The \(n\)-th time response of \(H\) and the spatial energy at height
\(n+\tfrac12\) carry the same scale weight.  This coincidence does not yet
relate their values.  The critical balance does.

Return to the exact balance already found in
\eqref{eq:BodySpatialEnergyBalanceIntro}:
\[
  E_s'=\mathcal P_s-2\nu E_{s+1}.
\]
At critical height,
\[
  H'=\mathcal P_{1/2}-2\nu E_{3/2}.
\]
Differentiate once and use the same balance at the next spatial height:
\[
\begin{aligned}
  H''
  &=\partial_t\mathcal P_{1/2}-2\nu E_{3/2}'
  \\
  &=\partial_t\mathcal P_{1/2}
  -2\nu\mathcal P_{3/2}
  +(2\nu)^2E_{5/2}.
\end{aligned}
\]
One more derivative gives
\[
  H'''
  =\partial_t^2\mathcal P_{1/2}
  -2\nu\partial_t\mathcal P_{3/2}
  +(2\nu)^2\mathcal P_{5/2}
  -(2\nu)^3E_{7/2}.
\]
The new top spatial energy has appeared, but nothing below it has disappeared.
Every differentiated transport--pressure response remains in the row.  The
pattern continues, and for \(n\geq1\), induction gives
\begin{equation}\label{eq:BodyCriticalHeightDiscovery}
  H^{(n)}
  =(-2\nu)^nE_{n+1/2}
  +\sum_{j=0}^{n-1}
  (-2\nu)^j
  \partial_t^{\,n-1-j}\mathcal P_{j+1/2}.
\end{equation}
Because \(\nu>0\), the completed row can be solved for its top spatial energy:
\begin{equation}\label{eq:BodyCompletedCriticalDiscovery}
  E_{n+1/2}
  =\frac{
    H^{(n)}-
    \displaystyle\sum_{j=0}^{n-1}
    (-2\nu)^j\partial_t^{\,n-1-j}\mathcal P_{j+1/2}
  }{(-2\nu)^n}.
\end{equation}
This algebra exposes the higher energy; it does not bound it.  Every transfer
term on the right remains, and controlling the new energy requires controlling
that entire lower triangular chain.  What has changed is the direction of the
problem.  Higher temporal response places
\[
  E_{1/2},
  \quad E_{3/2},
  \quad E_{5/2},
  \quad E_{7/2},
  \quad\ldots
\]
inside the completed critical-height Tower rather than carrying the argument
downward toward weaker spatial information.

The scalar rows still have to be tied back to the velocity that produced them.
Because
\[
  H(t)=\frac12
  \pair{\Lambda^{1/2}V_0(t)}
       {\Lambda^{1/2}V_0(t)}_{L^2},
\]
Leibniz's rule gives
\begin{equation}\label{eq:BodyHeightDerivativeTemporalContraction}
  H^{(n)}(t)
  =\frac12\sum_{a=0}^{n}\binom{n}{a}
  \operatorname{Re}
  \pair{\Lambda^{1/2}V_a(t)}
       {\Lambda^{1/2}V_{n-a}(t)}_{L^2}.
\end{equation}
Thus \(H^{(n)}\) is neither the response \(V_n\) nor the spatial energy
\(E_{n+1/2}\).  It is a quadratic contraction of the temporal responses, and
\eqref{eq:BodyCompletedCriticalDiscovery} can expose the top spatial energy
only after the complete transport--pressure row has been carried with it.

Only now do we have two independent directions that need one notation.  For a
temporal height \(n\) and a spatial multi-index \(\alpha\), set
\begin{equation}\label{eq:BodyMixedJetDefinition}
  J_{n,\alpha}:=\partial_t^n\partial_x^\alpha u,
  \qquad
  \Pi_{n,\alpha}:=\partial_t^n\partial_x^\alpha p,
  \qquad
  \alpha\in\Nzero^3.
\end{equation}
Differentiating the same momentum equation in both directions gives
\begin{align}
  J_{n+1,\alpha}
  &+
  \sum_{a=0}^{n}
  \sum_{\beta\leq\alpha}
  \binom{n}{a}\binom{\alpha}{\beta}
  (J_{a,\beta}\cdot\nabla)J_{n-a,\alpha-\beta}
  +\nabla\Pi_{n,\alpha}
  \notag\\
  &=\nu\Delta J_{n,\alpha},
  \qquad
  \nabla\cdot J_{n,\alpha}=0.
  \label{eq:BodyMixedJetRecurrence}
\end{align}
The pressure coordinate is still slaved to those velocity coordinates:
\begin{equation}\label{eq:BodyMixedPressureRecurrence}
  -\Delta\Pi_{n,\alpha}
  =\partial_x^\alpha\partial_i\partial_j
  \sum_{a=0}^{n}\binom{n}{a}(V_a)_i(V_{n-a})_j.
\end{equation}
The indices have not split the fluid into separate models.  They tell us which
time response of the same velocity is being read, and how finely that response
is being resolved in space.

Fix one temporal response \(V_r\).  Pairing its differentiated equation at a
spatial height \(s\) repeats the same adjacent movement.  Set
\[
  E_{r,s}:=\frac12\norm{\Lambda^sV_r}_2^2
\]
and let \(\mathcal P_{r,s}\) be the complete pairing of \(V_r\) with the
transport and pressure terms in its differentiated equation.  Then
\begin{equation}\label{eq:BodyEveryResponseDiscovery}
  E_{r,s}'=\mathcal P_{r,s}-2\nu E_{r,s+1}.
\end{equation}
The equation now names the spatial column that belongs to this response:
\[
  V_r\in H_x^0,
  \quad V_r\in H_x^1,
  \quad V_r\in H_x^2,
  \quad V_r\in H_x^3,
  \quad\ldots
\]
Temporal order chooses which response of the fluid is being read; spatial
order chooses how finely that response is resolved.  These membership
statements are the coordinates the proof has to establish, not consequences
of \eqref{eq:BodyEveryResponseDiscovery} by itself.  The nonlinear pairings in
\(\mathcal P_{r,s}\) still have to be bounded, first at a finite collection of
adjacent heights and eventually uniformly as the terminal time is approached.

This tells us why an intersection will eventually enter, but the order matters.
For a fixed finite \(r\) and \(m\), one must first prove that the corresponding
coordinate remains finite on the interval under consideration.  If that can be
done for every finite spatial height of one response, then
\[
  \bigcap_{m=0}^{\infty}H^m(\R^3)
  =H^\infty(\R^3)
  \hookrightarrow C^\infty(\R^3).
\]
If the same payment can then be made for every finite temporal response, the
common object is
\begin{equation}\label{eq:BodyDesiredMixedIntersection}
  \bigcap_{r=0}^{\infty}
  \bigcap_{m=0}^{\infty}
  H^r(I;H^m(\R^3)).
\end{equation}
Every prescribed temporal derivative and spatial derivative uses only a finite
coordinate, followed by the corresponding Sobolev embedding.  The infinite
intersection says that all of those finite payments belong to one compatible
history.  It cannot be inferred from the diagram alone.

We have reached the reversal without pretending to have completed it.  A
finite singular history must ascend through faster temporal responses while
losing spatial control.  Viscosity places an adjacent higher spatial energy in
every one of those response rows.  The remaining problem is to turn that
algebraic adjacency into estimates on finite pieces of the actual history and
then ask whether those estimates survive the terminal limit.  That finite
payment is what the two-dimensional rectangles are built to measure.
